\begin{table*}[htbp]
\centering
\caption{Requirement perturbations used in the preliminary study. These perturbations are adapted from prior studies on natural-language perturbation and prompt robustness.}
\label{tab:uncertainty_factors}
\begin{tabular*}{\textwidth}{@{\extracolsep{\fill}}p{3.3cm}p{8.0cm}p{3.0cm}p{2.0cm}@{}}
\toprule
\textbf{Perturbation Type} & \textbf{Explanation} & \textbf{Transformation Method} & \textbf{Reference} \\
\midrule

Content Rewriting
& Rewrite the requirement by abstracting or paraphrasing detailed expressions of functional constraints, edge cases, or task intent while keeping the same high-level programming task.
& LLM-assisted rewriting
& ~\cite{chen2026nlperturbator,zhu2024promptrobust} \\

Format Transformation
& Rewrite the requirement into a natural, conversational, and human-like paragraph, avoiding formal benchmark-style structures while preserving all functional requirements and constraints.
& LLM-assisted rewriting
& ~\cite{chen2026nlperturbator} \\

Example Removal
& Remove input-output examples from the requirement, such as doctest-style example usages in HumanEval and sample input/output cases in programming-contest-style tasks such as LiveCodeBench.
& Rule-based removal
& ~\cite{nashid2023retrieval,khojah2025impact} \\

Verbose Rewriting
& Add redundant descriptions or expand the requirement expression while preserving the intended functionality.
& Rule-based transformation
& ~\cite{chen2026nlperturbator,zhu2024promptrobust} \\

Noise Injection
& Insert irrelevant or distracting natural-language content into the requirement without introducing new inputs, outputs, constraints, or expected behavior.
& Rule-based transformation
& ~\cite{chen2026nlperturbator,zhu2024promptrobust} \\

Typographical Perturbation
& Introduce spelling or character-level errors while aiming to preserve the original requirement intent.
& TextAttack-based character transformation
& ~\cite{morris2020textattack,zhu2024promptrobust} \\

\bottomrule
\end{tabular*}
\end{table*}